\documentclass[border=10pt,tikz]{standalone}
\usepackage{fontspec}
\setmainfont{Apple SD Gothic Neo}
\setsansfont{Apple SD Gothic Neo}
\usepackage{tikz}
\usetikzlibrary{positioning,arrows.meta,shapes,fit,backgrounds,calc,decorations.pathreplacing}
\begin{document}
% _design.tex — shared TikZ design system for all blog diagrams
% Inputs nothing on its own — meant to be \input{} from each diagram .tex
% AFTER \begin{document}, BEFORE \begin{tikzpicture}.
%
% Provides a unified visual language across GoF / DSA / EMC++ / Beautiful series:
%   - Colors: muted, accessible, color-coded by role
%   - Node styles: consistent sizing, rounded, subtle borders
%   - Edge styles: UML-style inheritance, aggregation, plus dashed alternates
%
% Usage:
%   \documentclass[border=10pt,tikz]{standalone}
%   \usepackage{tikz}
%   \usetikzlibrary{positioning,arrows.meta,shapes,fit,backgrounds}
%   \begin{document}
%   \input{../_design.tex}     % adjust relative path
%   \begin{tikzpicture}[blog]  % apply 'blog' style
%   ...

\definecolor{absfill}{RGB}{220,232,247}   % cool blue
\definecolor{absbord}{RGB}{72,118,182}
\definecolor{confill}{RGB}{249,222,222}   % warm red
\definecolor{conbord}{RGB}{197,93,93}
\definecolor{impfill}{RGB}{223,238,222}   % soft green
\definecolor{impbord}{RGB}{99,154,93}
\definecolor{clifill}{RGB}{251,239,205}   % pale yellow
\definecolor{clibord}{RGB}{200,167,72}
\definecolor{hifill} {RGB}{253,224,200}   % accent orange
\definecolor{hibord} {RGB}{210,135,68}
\definecolor{nodefill}{RGB}{250,250,250}
\definecolor{nodebord}{RGB}{120,120,120}
\definecolor{edgec}  {RGB}{80,80,80}
\definecolor{edgesc} {RGB}{145,145,145}
\definecolor{groupbord}{RGB}{200,200,200}

\tikzset{
  blog/.style={
    font=\sffamily\small,
    every node/.style={font=\sffamily\small}
  },
  % --- node roles ---
  cls/.style={
    rectangle, rounded corners=2pt,
    draw=nodebord, line width=0.4pt,
    minimum height=8mm, minimum width=28mm,
    inner sep=3pt, fill=nodefill, align=center
  },
  abs/.style={cls, draw=absbord, fill=absfill, font=\sffamily\itshape\small},
  con/.style={cls, draw=conbord, fill=confill},
  imp/.style={cls, draw=impbord, fill=impfill},
  cli/.style={cls, draw=clibord, fill=clifill},
  hi/.style={cls, draw=hibord, fill=hifill, font=\sffamily\bfseries\small},
  % --- circle (for trees / graph nodes) ---
  vx/.style={
    circle, draw=absbord, line width=0.4pt,
    minimum size=8mm, fill=absfill, inner sep=0pt
  },
  vxc/.style={vx, draw=conbord, fill=confill},
  vxh/.style={vx, draw=hibord, fill=hifill, font=\sffamily\bfseries},
  % --- decision diamond ---
  q/.style={
    diamond, draw=clibord, line width=0.4pt, fill=clifill,
    inner sep=2pt, aspect=2.4, align=center
  },
  % --- edges ---
  inh/.style={                                  % UML inheritance
    -{Triangle[length=3mm,width=3mm,open]},
    draw=edgec, line width=0.5pt
  },
  agg/.style={                                  % UML aggregation
    {Diamond[length=2.5mm,width=2.5mm,open]}-{Stealth[length=2.2mm]},
    draw=edgec, line width=0.45pt
  },
  uses/.style={                                 % directed reference
    -{Stealth[length=2.2mm]}, draw=edgec, line width=0.5pt
  },
  arr/.style={                                  % plain arrow (graphs/trees)
    -{Stealth[length=2mm]}, draw=edgec, line width=0.45pt
  },
  edge/.style={                                 % undirected (trees)
    draw=edgec, line width=0.45pt
  },
  alt/.style={                                  % dashed: similar/alt
    -{Stealth[length=2mm]}, dashed, draw=edgesc, line width=0.45pt
  },
  ret/.style={                                  % dashed creates/returns
    -{Stealth[length=2mm]}, dashed, draw=edgesc, line width=0.45pt
  },
  thr/.style={                                  % threading link (red dashed)
    ->, dashed, draw=conbord, line width=0.5pt
  },
  % --- group / region ---
  group/.style={
    rectangle, rounded corners=4pt,
    draw=groupbord, line width=0.4pt,
    inner sep=10pt, fill=black!2
  },
  glabel/.style={font=\sffamily\bfseries\small, anchor=north west, inner sep=2pt}
}

% Korean font convenience macro — included by every Hangul diagram
\providecommand{\KOR}{}

\begin{tikzpicture}[blog, flowb/.style={cls, text width=5.6cm, minimum height=12mm}]
\node[font=\sffamily\bfseries\large] at (2.5,1.4) {Scan dump의 흐름};
\node[flowb, prp] (s0) at (0,0.0) {\textbf{1. 실패 순간을 잡는 트리거}\\[2pt]\footnotesize watchpoint·카운터·외부 신호};
\node[flowb, warn] (s1) at (0,-1.75) {\textbf{2. 클럭 정지}\\[2pt]\footnotesize 모든 플립플롭 값이 그대로 얼어붙음};
\draw[arr] (s0.south) -- (s1.north);
\node[flowb, hi] (s2) at (0,-3.5) {\textbf{3. scan 모드 진입}\\[2pt]\footnotesize SE=1, 플립플롭이 한 줄로 연결됨};
\draw[arr] (s1.south) -- (s2.north);
\node[flowb, hi] (s3) at (0,-5.25) {\textbf{4. shift out}\\[2pt]\footnotesize JTAG이나 테스터로 비트열을 끝까지 읽음};
\draw[arr] (s2.south) -- (s3.north);
\node[flowb, con] (s4) at (0,-7.0) {\textbf{5. 매핑 파일로 복원}\\[2pt]\footnotesize 비트 위치 → RTL 신호 이름};
\draw[arr] (s3.south) -- (s4.north);
\node[flowb, imp] (s5) at (0,-8.75) {\textbf{6. 시뮬레이션과 비교}\\[2pt]\footnotesize 같은 cycle의 기대 상태와 diff};
\draw[arr] (s4.south) -- (s5.north);
\node[note, text width=4.6cm, anchor=west] at (3.6,{-3*1.75}) {이 단계 이후 칩의 기능 상태는\\[2pt]망가집니다. 덤프가 끝나면\\[2pt]반드시 리셋해야 합니다.};
\draw[alt] (3.6,{-3*1.75}) -- (s3.east);
\node[note, text width=4.6cm, anchor=west] at (3.6,{-4*1.75-0.2}) {매핑 파일은 DFT 도구가\\[2pt]scan 삽입 때 만듭니다.};
\draw[alt] (3.6,{-4*1.75-0.2}) -- (s4.east);
\end{tikzpicture}
\end{document}
